\begin{PacketTable}{VCI Trust Anchor Status Reply (osdp\_PIV\_VCITASTATUSR)}
  \PacketRow{1}{CMND}{Reply identifier.}{0x88 (tentative)}
  \PacketRow{2}{TOTAL}{Length of the complete status payload, least-significant byte first.}{0x0000--0xFFFF}
  \PacketRow{2}{OFFSET}{Offset of this fragment within the status payload, least-significant byte first.}{0x0000--0xFFFF}
  \PacketRow{2}{DATA\_LEN}{Length of the data block in this fragment, least-significant byte first.}{0x0000--0xFFFF}
  \PacketRow{1}{Total Anchor Count}{Number of cached trust anchors.}{0x00--0xFF}
  \PacketRow{2}{Available Bytes}{Remaining storage capacity for additional variable-length trust anchor records, least-significant byte first.}{0x0000--0xFFFF}
  \PacketRow{0--n}{Anchor Records}{Zero or more 12-byte inventory records. Each record contains: \texttt{Trust Anchor ID} (1 byte for load/delete management), anchor selector \texttt{IIN} (8 bytes from the stored trust anchor record, not necessarily the secure CVC IIN), \texttt{Anchor Length} (2 bytes LSB first, maximum \texttt{0x0240}), and \texttt{Attributes} (1 byte, reserved). Records that span fragments continue in subsequent replies.}{--}
\end{PacketTable}
